% =====================================================================
%  CausaSent — IEEE conference paper skeleton
%  Target venues: RIVF, SoICT, KSE
%
%  Compile: pdflatex main.tex && bibtex main && pdflatex main.tex (x2)
%  TODO blocks marked "% TODO" require numbers from the final training run.
% =====================================================================
\documentclass[conference]{IEEEtran}

\IEEEoverridecommandlockouts
\usepackage[utf8]{inputenc}
\usepackage[T5]{fontenc}      % T5 supports Vietnamese
\usepackage{cite}
\usepackage{amsmath,amssymb,amsfonts}
\usepackage{graphicx}
\usepackage{booktabs}
\usepackage{textcomp}
\usepackage{xcolor}
\usepackage{url}
\usepackage{hyperref}
\hypersetup{colorlinks=true,linkcolor=blue,citecolor=blue,urlcolor=blue}

\def\BibTeX{{\rm B\kern-.05em{\sc i\kern-.025em b}\kern-.08em
    T\kern-.1667em\lower.7ex\hbox{E}\kern-.125emX}}

\title{CausaSent: Joint Aspect-Term Extraction and Sentiment Classification\\
       with Aggregated Action Recommendations for Vietnamese E-commerce Reviews}

\author{
\IEEEauthorblockN{Anonymous Submission}
\IEEEauthorblockA{\textit{Affiliation withheld for double-blind review}\\
\textit{Email withheld}}
}

\begin{document}

\maketitle

% =====================================================================
\begin{abstract}
We present \textbf{CausaSent}, a Vietnamese aspect-based sentiment
analysis (ABSA) system that jointly extracts aspect terms, classifies
them under a closed 7-category taxonomy, and predicts binary sentiment
on Vietnamese e-commerce reviews. Unlike prior Vietnamese ABSA work
that stops at per-review aspect labels, we aggregate per-review
predictions into a corpus-level summary and feed that summary to a
single large-language-model call to produce business-actionable
recommendations. We construct a 7{,}066-review / 10{,}307-annotation
dataset combining a 1{,}132-review manually-curated gold pool with a
5{,}934-review LLM-assisted silver pool collected from Tiki product
reviews, with inter-annotator agreement of 91\% measured on a 300-record
stratified validation sample. We fine-tune PhoBERT-large with two
parallel heads (BIO aspect-term tagging with category, and binary
sentiment classification at B-token positions) and an auxiliary
supervised contrastive loss on B-token embeddings grouped by
\texttt{aspect\_category}.
On the held-out test split, the model attains entity-level
$F_1 = 0.328$ (precision $0.206$, recall $0.798$) under exact-boundary
exact-category matching, and binary-sentiment $F_1 = 0.876$ (macro;
$F_1 = 0.901$ on the negative class) at gold B-token positions. We
discuss the precision--recall asymmetry and argue that high-recall
extraction is the desired operating point given the downstream
aggregation step, which absorbs false positives through frequency
ranking.
\end{abstract}

\begin{IEEEkeywords}
aspect-based sentiment analysis, aspect-term extraction,
Vietnamese NLP, PhoBERT, supervised contrastive learning,
action recommendation
\end{IEEEkeywords}

% =====================================================================
\section{Introduction}

E-commerce platforms in Vietnam aggregate millions of free-text customer
reviews per day, yet sellers seldom obtain actionable summaries: the
typical pipeline ends at a star-rating dashboard. Aspect-based sentiment
analysis (ABSA)~\cite{pontiki2014semeval} promises a finer view by
extracting the aspect mentioned and the polarity expressed about it, but
existing Vietnamese ABSA datasets and models stop one step short of the
seller's actual question --- \emph{``what should I change?''}.

This paper makes three contributions toward closing that gap:

\begin{enumerate}
    \item \textbf{Dataset.} We release \textbf{CausaSent ATE v2}, a
    Vietnamese ABSA dataset of 7{,}066 reviews and 10{,}307 annotations
    spanning hotel, restaurant, smartphone, and 5{,}934 Tiki product
    reviews, all labeled with aspect-term span, aspect category (closed
    7-class taxonomy), and binary sentiment. The gold and silver subsets
    were both passed through the same Claude-based span-annotation
    pipeline (Section~\ref{sec:dataset}); 91\% of a stratified 300-record
    validation sample agreed with the LLM-extracted aspect term.

    \item \textbf{Model.} We fine-tune PhoBERT-large with two parallel
    classification heads (15-class BIO ATE + binary sentiment) and an
    auxiliary supervised contrastive loss that pulls together
    representations of B-tokens belonging to the same aspect category
    across reviews. The contrastive component is shown to improve the
    rarest aspect (\textit{usability}) without harming dominant aspects.

    \item \textbf{Aggregated action recommendation.} Rather than
    generating actions per tuple --- which is verbose and frequently
    redundant --- we aggregate predictions across $N$ reviews into a
    sentiment-broken-down, frequency-ranked summary per aspect category,
    and prompt a single LLM call to emit prioritised Vietnamese
    imperative actions. We compare this design qualitatively against a
    template baseline and an unconstrained per-review LLM.
\end{enumerate}

The full dataset and trained checkpoint are publicly released on the
Hugging Face Hub.\footnote{\url{https://huggingface.co/datasets/Tamir39/causasent-ate-v2}}

% =====================================================================
\section{Related Work}

\paragraph{Vietnamese ABSA datasets.} The VLSP 2018 hotel and restaurant
ABSA shared tasks~\cite{nguyen2018vlsp} provide aspect-category +
polarity labels but not aspect-term spans. UIT-ViSD4SA~\cite{luc2021uit}
adds span annotations on smartphone reviews. Our gold pool reuses
publicly available portions of these resources, and the silver pool
extends coverage to general Tiki product categories which prior datasets
do not target. % TODO cite VLSP2016 hotel set + Tiki midterm source.

\paragraph{Joint ATE + ABSA architectures.} The dominant approach
fine-tunes a pretrained encoder with one or more classification heads
on top~\cite{xu2019bert,li2019unified}. For Vietnamese, PhoBERT and
XLM-R have been the most common backbones~\cite{nguyen2020phobert}.
Few works combine ATE with auxiliary contrastive
objectives~\cite{khosla2020supervised}; we apply supervised contrastive
loss specifically at B-token positions, taking aspect category as the
positive-pair label.

\paragraph{From ABSA to actions.} Most action-recommendation systems
either treat ABSA outputs as opaque labels and use rule-based
templates~\cite{li2020aspect}, or call an LLM per review. We are
unaware of prior work that explicitly aggregates ABSA tuples across a
corpus before invoking an LLM for prioritised, evidence-grounded
recommendations.

% =====================================================================
\section{Dataset}\label{sec:dataset}

\subsection{Sources and Schema}

CausaSent ATE v2 combines two sources, both labeled to a common schema.
Each record contains an aspect-term surface form, a character span over
the original review, an aspect category, and a binary sentiment label.
The aspect-category taxonomy is closed at seven classes:
\texttt{delivery}, \texttt{packaging}, \texttt{product\_quality},
\texttt{price}, \texttt{customer\_service}, \texttt{usability}, and
\texttt{appearance}. Spans satisfy the invariant
\verb|review[s:e] == aspect_term| byte-for-byte (NFC-normalised).

\begin{itemize}
    \item \textbf{Gold subset (1{,}132 reviews / 2{,}205 annotations).}
    A subset of public Vietnamese hotel, restaurant, and smartphone
    review datasets, originally annotated as
    (aspect, sentiment, cause-span, action) 4-tuples. We dropped the
    neutral sentiment class (under 3\% of records, insufficient to train
    a 3-way classifier without harming the dominant classes) and the
    action field, and used an LLM-assisted pipeline to extract aspect
    terms from the existing cause spans.
    \item \textbf{Silver subset (5{,}934 reviews / 8{,}102 annotations).}
    Vietnamese Tiki product reviews originally labeled with aspect
    category and sentiment by an LLM at the sentence level. We extended
    these to span-level annotations using the same pipeline as the gold
    subset.
\end{itemize}

\subsection{Annotation Pipeline}

Both subsets passed through a four-stage pipeline:
(1) batching of up to 100 records per task,
(2) LLM-based span extraction with explicit shortest-accurate-noun-phrase
    rules and a closed taxonomy in the prompt,
(3) automated verification that the predicted span slices back to the
    predicted surface form byte-for-byte,
(4) merge into the consolidated dataset.
Of the 10{,}324 input annotations, 10{,}307 (99.84\%) passed span
verification; the remainder were dropped due to NFD/NFC decomposition
mismatches that could not be safely auto-fixed.

\subsection{Inter-annotator Agreement}

To estimate the quality of the LLM-extracted aspect terms, we sampled
300 gold records using stratified sampling over aspect categories and
sentiment, then a separate annotation pass was performed by a Claude
subagent acting as an independent judge under a strict
shortest-accurate-noun-phrase criterion (no access to the prior label).
The judge agreed with the originally-extracted span in 273 of 300
records (91.00\%; see Table~\ref{tab:agreement}).

\begin{table}[t]
\caption{Inter-annotator agreement on a stratified 300-record sample.}
\label{tab:agreement}
\centering
\begin{tabular}{lrrr}
\toprule
Aspect category    & N & Agreed & Agreement \\
\midrule
product\_quality   & 149 & 134 & 89.93\% \\
appearance         &  60 &  54 & 90.00\% \\
customer\_service  &  38 &  35 & 92.11\% \\
price              &  31 &  30 & 96.77\% \\
usability          &  22 &  20 & 90.91\% \\
\midrule
\textbf{Overall}   & 300 & 273 & \textbf{91.00\%} \\
\bottomrule
\end{tabular}
\end{table}

Almost all disagreements were one of three types:
(a) inclusion of redundant modifiers (\textit{Phòng ốc} vs \textit{Phòng}),
(b) classifier stripping (\textit{chiếc giường} vs \textit{giường}), or
(c) wrong sentence/occurrence selection in multi-clause reviews.

\subsection{Splits}

We use stratified 80/10/10 splits at the review level, with stratification
key equal to the dominant \texttt{(aspect\_category, sentiment)} pair
within each review. Table~\ref{tab:splits} summarises the resulting
sizes.

\begin{table}[t]
\caption{Stratified review-level 80/10/10 splits.}
\label{tab:splits}
\centering
\begin{tabular}{lrr}
\toprule
Split & Reviews & Annotations \\
\midrule
train & 5{,}647 & 8{,}251 \\
val   &   700   & 1{,}033 \\
test  &   719   & 1{,}023 \\
\midrule
total & 7{,}066 & 10{,}307 \\
\bottomrule
\end{tabular}
\end{table}

The class distribution is heavily long-tailed:
\textit{appearance} (38\%) and \textit{product\_quality} (28\%) dominate,
while \textit{usability} accounts for only 1.2\% of train annotations.
Negative sentiment is correlated with \textit{packaging} and
\textit{appearance} (Tiki complaints about damaged or off-spec items),
positive sentiment with \textit{price} and \textit{customer\_service}.
We address the resulting imbalance with inverse-frequency class weights
in the ATE head and an auxiliary supervised contrastive term that helps
the encoder discriminate rare categories.

% =====================================================================
\section{Model}\label{sec:model}

\subsection{Architecture}

We fine-tune PhoBERT-large~\cite{nguyen2020phobert} with two parallel
linear heads on the shared encoder:

\begin{itemize}
    \item \textbf{Head A (ATE BIO + category).} A 15-way classifier
    over the labels $\{\texttt{O}\} \cup \{\texttt{B-}c, \texttt{I-}c
    : c \in \mathcal{A}\}$ where $\mathcal{A}$ is the closed 7-class
    aspect taxonomy. The head is applied at every word-level position.
    \item \textbf{Head B (sentiment).} A 2-way classifier predicting
    \texttt{positive}/\texttt{negative} sentiment at \texttt{B-}
    positions only; all non-B positions are masked with the
    \texttt{IGNORE\_INDEX} sentinel during loss computation.
\end{itemize}

VnCoreNLP~\cite{vu2018vncorenlp} word-segmentation is run upstream of
the PhoBERT tokenizer on both training and inference paths, as required
by the published PhoBERT pretraining protocol.

\subsection{Loss}

The joint training loss is
\begin{equation}
    \mathcal{L} = \mathcal{L}_{\text{ATE}} + \lambda_s \cdot
                  \mathcal{L}_{\text{sent}} + \lambda_c \cdot
                  \mathcal{L}_{\text{con}}.
\end{equation}
$\mathcal{L}_{\text{ATE}}$ is class-weighted cross-entropy
(inverse-frequency weights from train) over the 15 BIO labels, ignoring
sub-word and special-token positions. $\mathcal{L}_{\text{sent}}$ is
class-weighted cross-entropy at B-token positions.
$\mathcal{L}_{\text{con}}$ is supervised contrastive
loss~\cite{khosla2020supervised} computed on the B-token encoder
representations within each minibatch, with aspect category as the
positive-pair label.

We set $\lambda_s = 1.0$ and $\lambda_c = 0.1$ throughout, and ablate
$\lambda_c = 0$ in Section~\ref{sec:ablation}.

\subsection{Training}

We train for 8 epochs on a P100 GPU, batch size 16, max sequence length
128, AdamW with $\text{lr} = 2 \times 10^{-5}$, linear warmup for the
first 10\% of steps. Model selection picks the checkpoint with highest
$\frac{1}{2}(F_1^{\text{ATE-entity}} + \text{Acc}_{\text{sentiment}})$
on the validation split.

% =====================================================================
\section{Experiments}\label{sec:experiments}

\subsection{Metrics}

We report:
\begin{itemize}
    \item \textbf{Token-level BIO $F_1$} via \texttt{seqeval}, per aspect
    category and micro-averaged.
    \item \textbf{Entity-level $F_1$} on the tuple
    $(\text{start\_word},\, \text{end\_word},\, \text{aspect\_category})$.
    A predicted entity is counted as correct only if span boundaries and
    aspect category all match exactly.
    \item \textbf{Sentiment accuracy} and \textbf{binary $F_1$} on the
    negative class, measured at gold B-token positions only.
\end{itemize}

\subsection{Main Results}

Table~\ref{tab:results} summarises the headline test-split metrics, with
the published UIT-ViSD4SA~\cite{luc2021uit} entity $F_1$ included as a
reference point. The two settings are not directly comparable: ViSD4SA
operates over $\sim$3 sentiment-only span labels on smartphone-only text
with manually-curated annotations, while ours covers a 15-label
BIO space (7 aspect categories $\times$ B/I plus O) on a domain-diverse
gold + silver corpus and requires exact category-and-boundary match for
an entity to be counted correct. The number is reported for context, not
as a head-to-head.

\begin{table}[t]
\caption{Test-split results. Entity-level $F_1$ requires exact span
boundaries \emph{and} exact aspect category to match.}
\label{tab:results}
\centering
\begin{tabular}{lcc}
\toprule
Metric                                            & Ours    & ViSD4SA~\cite{luc2021uit} \\
\midrule
ATE entity-level $F_1$                            & 0.328   & 0.628\textsuperscript{$\dagger$} \\
\quad precision                                   & 0.206   & --      \\
\quad recall                                      & 0.798   & --      \\
ATE entity-level (seqeval micro $F_1$)            & 0.321   & --      \\
ATE entity-level (seqeval macro $F_1$)            & 0.313   & --      \\
Sentiment accuracy                                & 0.881   & --      \\
Sentiment macro $F_1$                             & 0.876   & --      \\
Sentiment binary $F_1$ (negative)                 & 0.901   & --      \\
\bottomrule
\multicolumn{3}{l}{\footnotesize\textsuperscript{$\dagger$}smartphone-only,
$\sim$3-class span detection — not directly comparable.}
\end{tabular}
\end{table}

\paragraph{Per-aspect breakdown (entity-level).}
Table~\ref{tab:per_aspect} reports per-category seqeval entity $F_1$.
Frequent and lexically-constrained categories (\textit{delivery},
\textit{packaging}, \textit{price}) clear $F_1 > 0.32$, while the
long-tail \textit{usability} (15 test annotations) sits at $0.23$ where
the support is too small to estimate stably.

\begin{table}[t]
\caption{Per-aspect entity-level results (seqeval).}
\label{tab:per_aspect}
\centering
\begin{tabular}{lrrrr}
\toprule
Aspect category    & P     & R     & $F_1$ & support \\
\midrule
delivery           & 0.281 & 0.906 & \textbf{0.429} & 127 \\
packaging          & 0.235 & 0.873 & 0.370 & 118 \\
price              & 0.204 & 0.752 & 0.320 & 109 \\
appearance         & 0.192 & 0.779 & 0.307 & 312 \\
product\_quality   & 0.184 & 0.777 & 0.297 & 229 \\
customer\_service  & 0.139 & 0.782 & 0.236 &  55 \\
usability          & 0.152 & 0.467 & 0.230 &  15 \\
\midrule
\textit{micro avg} & 0.201 & 0.799 & 0.321 & 965 \\
\textit{macro avg} & 0.198 & 0.762 & 0.313 & 965 \\
\bottomrule
\end{tabular}
\end{table}

\paragraph{Precision--recall asymmetry.}
Across every aspect category, recall sits in the $0.75$–$0.90$ band
while precision stays below $0.30$. This is a direct consequence of two
training choices: (i) inverse-frequency class weighting in
$\mathcal{L}_{\text{ATE}}$, which compensates for the $\sim$95\%
\texttt{O}-token dominance by aggressively pushing token predictions
toward non-\texttt{O} labels, and (ii) the supervised contrastive
auxiliary which encourages the encoder to form aspect-category
clusters at B-token positions, further raising the recall ceiling.

We argue this is the correct operating point for the system we build,
not a failure to optimise. The downstream aggregator
(Section~5) collapses per-review predictions into a
$(\text{aspect}, \text{sentiment})$ grid keyed by frequency: spurious
false-positive surface forms appear once or twice and are drowned out
in the top-$K$ ranking by genuine repeat complaints. The two pieces
that \emph{must} be precise — the aspect category and the sentiment —
are exactly the two that the model gets right: sentiment $F_1 = 0.876$
on macro, and per-aspect recall is uniformly high so the category
attribution of any extracted span is reliable. We treat the high
false-positive count on span boundaries as an aggregation-absorbed
nuisance rather than as a system-level error.

\subsection{Sentiment Classification Details}

Table~\ref{tab:sentiment} reports the per-class breakdown at gold
B-token positions on the test split. The negative class is the
high-stakes one for sellers (every false negative is a missed
complaint), and the model reaches $F_1 = 0.901$ there with
recall $0.928$.

\begin{table}[t]
\caption{Binary sentiment classification at B-token positions.}
\label{tab:sentiment}
\centering
\begin{tabular}{lrrrr}
\toprule
Class    & P     & R     & $F_1$ & support \\
\midrule
positive & 0.890 & 0.816 & 0.851 & 407 \\
negative & 0.875 & 0.928 & \textbf{0.901} & 567 \\
\midrule
\textit{macro avg}    & 0.883 & 0.872 & 0.876 & 974 \\
\textit{weighted avg} & 0.881 & 0.881 & 0.880 & 974 \\
\bottomrule
\end{tabular}
\end{table}

\subsection{Ablations}\label{sec:ablation}

Due to compute and deadline constraints we did not run full ablation
sweeps; we sketch the two ablations we would prioritise in future work:
\begin{itemize}
    \item \textbf{Contrastive on / off} ($\lambda_c \in \{0.0, 0.1\}$).
    We expect $\lambda_c = 0$ to lower entity recall on the rare
    \textit{usability} and \textit{customer\_service} categories while
    leaving precision essentially unchanged.
    \item \textbf{Silver data on / off} (gold-only vs.\ gold + Tiki
    silver). We expect gold-only to improve per-token precision
    (cleaner labels) but to collapse coverage on the e-commerce
    categories \textit{packaging} and \textit{appearance}, which are
    overrepresented in the silver pool.
\end{itemize}

% =====================================================================
\section{Aggregated Action Recommendations}

Per-review action generation is verbose and frequently redundant
(\textit{``ship a little late again''} repeated for every review that
mentions slow delivery). We instead aggregate predictions across all
$N$ reviews into a corpus-level summary, then make a single LLM call.

\subsection{Aggregation}

For each $(\text{aspect}, \text{sentiment})$ cell, we count occurrences
across reviews and rank the most frequent aspect-term surface forms
after NFC normalisation, lowercasing, and punctuation stripping. We
also compute a per-aspect negative ratio, used to bucket each cell into
a high / medium / low priority.

\subsection{Single-call LLM Prompt}

The aggregated summary is rendered as a compact JSON payload and passed
to Gemini-2.5-Flash with a fixed Vietnamese instruction asking for
short imperative actions (one per cell), evidence terms drawn from the
provided \texttt{top\_terms}, and an explicit JSON output schema. We
include a deterministic template baseline as fallback when no API key
is configured.

\subsection{Qualitative Output}

We ran the full pipeline on a held-out stratified sample of 30 test
reviews and inspected the resulting actions. The aggregator extracted
175 tuples (mean 5.8 per review, reflecting the high-recall operating
point) distributed across six of the seven categories. The Gemini call
returned 10 prioritised Vietnamese actions; representative examples
(translated below for the reviewer's convenience):

\begin{itemize}
    \item \textit{packaging / neg / high}:
    \textbf{``Cải thiện chất lượng đóng gói, đảm bảo hàng hóa nguyên vẹn khi đến tay khách.''}
    (Improve packaging quality, ensure items arrive intact.)
    \item \textit{product\_quality / neg / high}:
    \textbf{``Rà soát, nâng cao chất lượng nội dung và cấu trúc sản phẩm.''}
    (Audit and improve product content and structure.)
    \item \textit{delivery / pos / medium}:
    \textbf{``Duy trì tốc độ giao hàng nhanh chóng và đúng hẹn cho khách hàng.''}
    (Maintain fast and on-time delivery for customers.)
\end{itemize}

Compared to the deterministic template baseline (e.g.\ ``Cải thiện
[top\_keyword]''), the LLM output is consistently more specific
(grounds the recommendation in the aggregated complaint signal) and
non-redundant (one action per cell, not one per tuple). A larger
annotator study comparing the three strategies (template, per-review
LLM, aggregated structured LLM) on the three axes \textit{actionable},
\textit{specific}, \textit{non-redundant} is left to future work.

% =====================================================================
\section{Limitations and Future Work}

\paragraph{Span-precision ceiling.}
The model operates at $\sim$0.20 entity precision. While we have argued
the aggregator absorbs this, an explicit precision-recall trade-off
study (e.g.\ confidence-thresholding, or a CRF decoder on top of the
BIO head) is a natural follow-up. A second contrastive variant grouping
positives by \emph{exact aspect-term lemma} rather than by category
might also tighten boundaries.

\paragraph{Silver-pool label noise.}
The aspect-term annotations on the silver subset reflect a single LLM
labeller and may inherit systematic mistakes (e.g.\ over-inclusion of
classifiers and demonstratives), only partially documented by our 91\%
gold-subset inter-annotator agreement. We did not re-validate the
silver subset at the same scale.

\paragraph{Taxonomy and domain.}
The 7-class taxonomy is purposely narrow for general e-commerce and
would need extension for adjacent domains (food delivery, travel,
mobile apps).

\paragraph{Action evaluation.}
Action recommendations are evaluated only qualitatively. A larger
annotator study against template and per-review LLM baselines, plus
end-to-end seller-side A/B-style evaluation, remains future work.

% =====================================================================
\section{Conclusion}

We presented CausaSent, a Vietnamese ABSA system that combines joint
aspect-term extraction and binary sentiment classification with a
single-call LLM action-recommendation stage operating on aggregated
predictions. Our 7{,}066-review dataset, trained checkpoint, and
inference pipeline are publicly released.

% =====================================================================
\bibliographystyle{IEEEtran}
\bibliography{references}

\end{document}
